\documentclass{article}
\usepackage{amsmath,amssymb,amsthm}
\usepackage{bm}
\usepackage{algorithm}
\usepackage{algpseudocode}
\usepackage{hyperref}
\usepackage{enumitem}
\usepackage{geometry}
\geometry{margin=1in}
\newtheorem{theorem}{Theorem}
\newtheorem{lemma}{Lemma}
\newcommand{\E}{\mathbb{E}}
\newcommand{\R}{\mathbb{R}}
\newcommand{\norm}[1]{\left\lVert #1\right\rVert}
\newcommand{\ip}[2]{\left\langle #1,#2\right\rangle}
\newcommand{\cond}[1]{\quad\text{#1}}
\begin{document}

\section*{Algorithm Name}
\textbf{Stochastic Newton with Bounded Hessian Estimates (SN-BH)}  

The method updates the iterate $x_t\in\R^{d}$ by a pre‑conditioned gradient step in which the pre‑conditioner is a stochastic estimate $H_t^{-1}$ of the true Hessian $\nabla^{2}f(x_t)$.  

\section*{Mathematical Formulation}
Let $f:\R^{d}\rightarrow\R$ be \emph{strongly convex} and \emph{smooth}.  
For a given step size $\alpha>0$ the iteration reads
\begin{align}
   g_t &:= \nabla f(x_t) , \label{eq:grad}\\
   H_t &\in\R^{d\times d}\ \text{ stochastic Hessian estimate},\qquad 
          H_t \succ 0, \label{eq:hessian}\\
   x_{t+1} &:= x_t - \alpha\, H_t^{-1} g_t . \label{eq:update}
\end{align}
The algorithm requires only that $H_t$ be an unbiased, uniformly well‑conditioned
estimate of the true Hessian:
\[
   \E\!\left[\,H_t\mid x_t\right]=\nabla^{2}f(x_t),\qquad
   \mu I \preceq H_t \preceq L I\quad\text{a.s.} \tag{A3}
\]
where $0<\mu\le L$ are the strong‑convexity and smoothness constants of $f$.

\section*{Pseudocode}
\begin{algorithm}[H]
\caption{SN‑BH}
\begin{algorithmic}[1]
\Require Initial point $x_0\in\R^{d}$, step size $\alpha\in(0,2/\kappa)$,
        constants $\mu,L$ with $\kappa:=L/\mu$.
\For{$t=0,1,2,\dots,T-1$}
    \State Compute (or sample) stochastic gradient $g_t=\nabla f(x_t)$
    \State Obtain stochastic Hessian estimate $H_t$ satisfying (A3)
    \State $x_{t+1}\gets x_t-\alpha\,H_t^{-1}g_t$
\EndFor
\State \Return $x_T$
\end{algorithmic}
\end{algorithm}

\section*{Dry Run Example}
Consider the one‑dimensional quadratic
\[
   f(w)=(w-3)^2.
\]
We have $\nabla f(w)=2(w-3)$, $\nabla^{2}f(w)=2$, thus $\mu=L=2$ and $\kappa=1$.  
Choose $\alpha=0.1$, $w_0=0$, and use the exact Hessian $H_t=2$ (which trivially satisfies (A3)).  

\begin{center}
\begin{tabular}{c|c|c|c}
$t$ & $w_t$ & $\nabla f(w_t)$ & $f(w_t)$\\ \hline
0 & $0$ & $-6$ & $9$\\
1 & $0.3$ & $-5.4$ & $(0.3-3)^2=7.29$\\
2 & $0.57$ & $-4.86$ & $(0.57-3)^2=5.9049$\\
3 & $0.813$ & $-4.374$ & $(0.813-3)^2=4.7449$
\end{tabular}
\end{center}
The update rule \eqref{eq:update} with $H_t^{-1}=1/2$ gives
\[
   w_{t+1}=w_t-\alpha\frac{1}{2}\,\nabla f(w_t)
          =w_t+0.1\cdot\frac{1}{2}\cdot(-\nabla f(w_t)),
\]
which reproduces the numbers above and shows a steady linear decrease of the objective.

\newpage
\section*{Assumptions}
\begin{enumerate}[label=(A\arabic*)]
\item \textbf{Strong convexity.} $f$ satisfies
\[
   f(y)\ge f(x)+\ip{\nabla f(x)}{y-x}+\frac{\mu}{2}\norm{y-x}^{2},
   \quad\forall x,y\in\R^{d}. \tag{A1}
\]
\item \textbf{Smoothness.} $f$ has $L$‑Lipschitz gradients:
\[
   \norm{\nabla f(y)-\nabla f(x)}\le L\norm{y-x},
   \quad\forall x,y\in\R^{d}. \tag{A2}
\]
\item \textbf{Unbiased, uniformly well‑conditioned Hessian estimates.}
\[
   \E\!\left[H_t\mid x_t\right]=\nabla^{2}f(x_t),\qquad
   \mu I\preceq H_t\preceq L I\quad\text{a.s.} \tag{A3}
\]
\item \textbf{Exact gradient.} For clarity of exposition we assume that the
gradient used in \eqref{eq:grad} is exact (the extension to unbiased stochastic
gradients follows by the same steps, adding a variance term).
\end{enumerate}

Define the condition number $\kappa:=L/\mu\ge1$.

\section*{Main Convergence Theorem}
\begin{theorem}[Linear convergence of SN‑BH]\label{thm:linear}
Assume (A1)–(A3) and fix a stepsize $\alpha\in\bigl(0,\frac{2}{\kappa}\bigr)$.  
Let $x^{\star}=\arg\min_{x}f(x)$ be the unique minimiser. Then for all $t\ge0$
\begin{equation}
   \E\!\left[\norm{x_{t+1}-x^{\star}}^{2}\mid x_t\right]
   \;\le\;
   \bigl(1-\rho\bigr)\,\norm{x_t-x^{\star}}^{2},
   \qquad\text{where }\;
   \rho:=\frac{2\alpha}{\kappa}-\alpha^{2}\kappa \in(0,1). \tag{T1}
\end{equation}
Consequently,
\begin{equation}
   \E\!\bigl[\norm{x_{t}-x^{\star}}^{2}\bigr]
   \le(1-\rho)^{t}\,\norm{x_{0}-x^{\star}}^{2},
   \qquad\forall t\ge0. \tag{T2}
\end{equation}
Thus SN‑BH enjoys a global linear (geometric) convergence rate.
\end{theorem}

\section*{Theoretical Properties}
\begin{enumerate}[label=\arabic*.]
\item \textbf{Stability.} The contraction factor $\rho$ is positive precisely when
$\alpha\in(0,2/\kappa)$. Choosing $\alpha$ close to the upper bound yields the
fastest possible linear rate $1-\rho\approx 1-1/\kappa$.
\item \textbf{Hyperparameter sensitivity.} The only tunable scalar is $\alpha$;
its admissible interval depends solely on the condition number $\kappa$,
which can be estimated from a few Hessian samples.
\item \textbf{Comparison with baselines.}
\begin{itemize}
    \item Gradient descent with step size $1/L$ contracts by a factor $1-\mu/L=1-1/\kappa$.
    \item SN‑BH contracts by $1-\rho$, where $\rho\ge\frac{1}{\kappa}$ for the optimal
          $\alpha=\frac{1}{\kappa}$, giving a strictly better rate.
    \item When $H_t=\nabla^{2}f(x_t)$ exactly and $\alpha=1$, the method reduces to the classic Newton step and enjoys a \emph{quadratic} local convergence; the theorem recovers the global linear regime that precedes the quadratic phase.
\end{itemize}
\end{enumerate}

\section*{Discussion}
The theorem shows that as long as the stochastic Hessian estimates remain uniformly
well‑conditioned (A3), the algorithm behaves like a pre‑conditioned gradient descent
with a fixed positive definite matrix bounded between $\mu I$ and $L I$.  The
expectation operator eliminates the randomness of $H_t$ while preserving the
contraction property because the bounds hold \emph{almost surely}.  Extensions to
unbiased stochastic gradients, mini‑batching, or adaptive stepsizes follow by
adding standard variance terms to the recursion \eqref{T1}.

\newpage
\section*{Preliminaries (Definitions and Lemmas)}
\begin{lemma}[Bounds on the inverse Hessian]\label{lem:inv}
If $\mu I\preceq H\preceq L I$, then
\[
   \frac{1}{L} I\preceq H^{-1}\preceq \frac{1}{\mu} I .
\]
\end{lemma}
\begin{proof}
Since $H\succeq\mu I$, for any $v\neq0$,
\[
   v^{\!\top}Hv\ge\mu\norm{v}^{2}\;\Longrightarrow\;
   v^{\!\top}H^{-1}v\le\frac{1}{\mu}\norm{v}^{2}.
\]
Similarly $H\preceq L I$ implies $v^{\!\top}H^{-1}v\ge\frac{1}{L}\norm{v}^{2}$.
\end{proof}

\begin{lemma}[Strong convexity $\Rightarrow$ gradient lower bound]\label{lem:sc}
For a $\mu$‑strongly convex $f$, the following holds for all $x$:
\[
   \ip{\nabla f(x)}{x-x^{\star}}\ge\mu\norm{x-x^{\star}}^{2}. \tag{L1}
\]
\end{lemma}
\begin{proof}
Apply (A1) with $y=x^{\star}$ and use $\nabla f(x^{\star})=0$:
\[
   f(x^{\star})\ge f(x)+\ip{\nabla f(x)}{x^{\star}-x}+\frac{\mu}{2}\norm{x^{\star}-x}^{2}.
\]
Rearranging gives (L1).
\end{proof}

\begin{lemma}[Smoothness $\Rightarrow$ gradient upper bound]\label{lem:smooth}
If $\nabla f$ is $L$‑Lipschitz, then for all $x$,
\[
   \norm{\nabla f(x)}\le L\norm{x-x^{\star}}. \tag{L2}
\]
\end{lemma}
\begin{proof}
From $L$‑Lipschitzness,
\[
   \norm{\nabla f(x)-\nabla f(x^{\star})}\le L\norm{x-x^{\star}}.
\]
Since $\nabla f(x^{\star})=0$, (L2) follows.
\end{proof}

\section*{Main Proof (Step‑by‑Step Derivation)}
We prove Theorem~\ref{thm:linear}.  All non‑trivial steps are numbered for easy reference.

\begin{proof}
Let $e_t:=x_t-x^{\star}$.  Using the update \eqref{eq:update} we have
\begin{align}
   e_{t+1}
   &= e_t - \alpha H_t^{-1}\nabla f(x_t). \tag{Step 1}
\end{align}
Taking the squared Euclidean norm and expanding,
\begin{align}
   \norm{e_{t+1}}^{2}
   &= \norm{e_t}^{2}
      -2\alpha\ip{H_t^{-1}\nabla f(x_t)}{e_t}
      +\alpha^{2}\norm{H_t^{-1}\nabla f(x_t)}^{2}. \tag{Step 2}
\end{align}
Condition on $x_t$ and use Lemma~\ref{lem:inv} to bound the two inner products.

\medskip\noindent\textbf{Bounding the cross term.}
Since $H_t^{-1}\preceq\frac{1}{\mu}I$ and $H_t^{-1}\succeq\frac{1}{L}I$, we obtain
\begin{align}
   \ip{H_t^{-1}\nabla f(x_t)}{e_t}
   &\ge \frac{1}{L}\ip{\nabla f(x_t)}{e_t} \tag{Step 3}\\
   &\ge \frac{\mu}{L}\norm{e_t}^{2} \qquad\text{(by Lemma~\ref{lem:sc})}. \tag{Step 4}
\end{align}
The first inequality uses the lower bound on $H_t^{-1}$; the second uses (L1).

\medskip\noindent\textbf{Bounding the quadratic term.}
Similarly,
\begin{align}
   \norm{H_t^{-1}\nabla f(x_t)}^{2}
   &\le \frac{1}{\mu^{2}}\norm{\nabla f(x_t)}^{2}
        \qquad\text{(by Lemma~\ref{lem:inv})}. \tag{Step 5}\\
   &\le \frac{L^{2}}{\mu^{2}}\norm{e_t}^{2}
        \qquad\text{(by Lemma~\ref{lem:smooth})}. \tag{Step 6}
\end{align}
Insert the bounds \eqref{Step 4} and \eqref{Step 6} into \eqref{Step 2}:
\begin{align}
   \norm{e_{t+1}}^{2}
   &\le \norm{e_t}^{2}
        -2\alpha\frac{\mu}{L}\norm{e_t}^{2}
        +\alpha^{2}\frac{L^{2}}{\mu^{2}}\norm{e_t}^{2}. \tag{Step 7}
\end{align}
Factor $\norm{e_t}^{2}$ and introduce the condition number $\kappa:=L/\mu$:
\begin{align}
   \norm{e_{t+1}}^{2}
   &\le \Bigl(1-\frac{2\alpha}{\kappa}+\alpha^{2}\kappa^{2}\Bigr)\norm{e_t}^{2}. \tag{Step 8}
\end{align}
Define the contraction coefficient
\[
   q(\alpha):=1-\frac{2\alpha}{\kappa}+\alpha^{2}\kappa^{2}.
\]
Since $q(\alpha)$ is a convex quadratic in $\alpha$, it attains its minimum at
$\alpha^{\star}=1/\kappa^{2}$.  For any $\alpha\in(0,2/\kappa)$ we have
$0<q(\alpha)<1$.  Set
\[
   \rho:=1-q(\alpha)=\frac{2\alpha}{\kappa}-\alpha^{2}\kappa
        \in(0,1). \tag{Step 9}
\]
Taking expectation conditioned on $x_t$ and using that the bound \eqref{Step 8}
holds almost surely (by (A3)), we obtain
\begin{align}
   \E\!\bigl[\norm{e_{t+1}}^{2}\mid x_t\bigr]
   &\le (1-\rho)\,\norm{e_t}^{2}. \tag{Step 10}
\end{align}
Unconditioning yields the recursion
\[
   \E\!\bigl[\norm{e_{t+1}}^{2}\bigr]
   \le (1-\rho)\,\E\!\bigl[\norm{e_t}^{2}\bigr].
\]
Iterating this inequality $t$ times gives
\[
   \E\!\bigl[\norm{e_{t}}^{2}\bigr]
   \le (1-\rho)^{t}\,\norm{e_{0}}^{2}, \tag{Step 11}
\]
which is exactly \eqref{T2}.  This completes the proof. \qed
\end{proof}

\section*{Rate Derivation (With Constants)}
From the definition of $\rho$ in \eqref{Step 9},
\[
   \rho(\alpha)=\frac{2\alpha}{\kappa}-\alpha^{2}\kappa.
\]
The optimal constant stepsize maximising $\rho$ on $(0,2/\kappa)$ solves
$\partial_{\alpha}\rho=0$, i.e.
\[
   \frac{2}{\kappa}-2\alpha\kappa=0
   \;\Longrightarrow\;
   \alpha_{\text{opt}}=\frac{1}{\kappa^{2}}.
\]
Substituting $\alpha_{\text{opt}}$ yields the best achievable contraction
\[
   \rho_{\max}= \frac{1}{\kappa^{2}}.
\]
Hence the linear rate can be expressed as
\[
   \E\!\bigl[\norm{x_{t}-x^{\star}}^{2}\bigr]
   \le\Bigl(1-\frac{1}{\kappa^{2}}\Bigr)^{t}\norm{x_{0}-x^{\star}}^{2}.
\]
If a more aggressive stepsize $\alpha=1/\kappa$ is used (still admissible for
$\kappa\ge2$), then
\[
   \rho = \frac{2}{\kappa^{2}}-\frac{1}{\kappa}= \frac{2-\kappa}{\kappa^{2}}<\frac{1}{\kappa^{2}},
\]
illustrating the trade‑off between step size and contraction.

\section*{Special Cases}
\begin{enumerate}[label=\alph*)]
\item \textbf{Exact Hessian.} If $H_t=\nabla^{2}f(x_t)$ deterministically,
    Lemma~\ref{lem:inv} still holds, and the same linear bound applies.
    In a neighbourhood where $\nabla^{2}f$ is Lipschitz, the recursion
    tightens to the classic Newton quadratic convergence.
\item \textbf{Stochastic gradient.} If $g_t$ is an unbiased estimator of
    $\nabla f(x_t)$ with bounded variance $\sigma^{2}$,
    an extra term $\alpha^{2}\sigma^{2}/\mu^{2}$ appears on the right‑hand side of
    \eqref{Step 7}.  The recursion becomes
    \[
       \E\!\bigl[\norm{e_{t+1}}^{2}\bigr]
       \le (1-\rho)\E\!\bigl[\norm{e_t}^{2}\bigr]
          +\frac{\alpha^{2}\sigma^{2}}{\mu^{2}}.
    \]
    Consequently the iterates converge linearly to a neighbourhood of $x^{\star}$
    whose radius is $O(\alpha\sigma/\mu)$.
\item \textbf{Mini‑batch Hessian.} Averaging $B$ independent Hessian samples
    reduces the variance of $H_t$ and thus relaxes the bound (A3) to
    $\mu I\preceq \E[H_t\mid x_t]\preceq L I$ while the almost‑sure bounds improve
    proportionally to $1/B$.
\end{enumerate}

\end{document}